\documentclass[10pt,a4paper]{article}

% =========================================================
% PACKAGES
% =========================================================
\usepackage[a4paper,margin=0.55in]{geometry}
\usepackage[fontsize=10.5pt]{scrextend}
\usepackage{times}
\usepackage{amsmath,amssymb}
\usepackage{booktabs}
\usepackage{array}
\usepackage{enumitem}
\usepackage{titlesec}
\usepackage{microtype}
\usepackage{balance}
\usepackage{graphicx}

% =========================================================
% TWO-COLUMN LAYOUT
% =========================================================
\setlength{\columnsep}{0.18in}
\setlength{\parindent}{0pt}
\setlength{\parskip}{2.2pt}

% =========================================================
% TYPOGRAPHY
% =========================================================
\titleformat{\section}
  {\bfseries\normalsize}
  {\Roman{section}.}
  {0.35em}
  {}

\titlespacing*{\section}{0pt}{5pt}{2pt}

\titleformat{\subsection}
  {\bfseries\small}
  {}
  {0pt}
  {}

\titlespacing*{\subsection}{0pt}{3pt}{1pt}

\setlist[itemize]{
  leftmargin=1.2em,
  itemsep=0pt,
  topsep=1pt,
  parsep=0pt
}

\setlist[enumerate]{
  leftmargin=1.4em,
  itemsep=0pt,
  topsep=1pt,
  parsep=0pt
}

% Slightly more readable tables
\renewcommand{\arraystretch}{1.05}

% =========================================================
% DOCUMENT
% =========================================================
\begin{document}

% =========================================================
% TITLE PAGE
% =========================================================
\begin{titlepage}
\begin{center}
\vspace*{0.5in}

{\large\bfseries BTP-I Mid-Semester Project Summary}

\vspace{0.4in}

{\Large\bfseries Dynamic Truck--Drone Collaborative Delivery\\[2pt] with Flexible Launch Positions}

\vspace{0.55in}

by

\vspace{0.12in}

{\large Jeet Anand}\\
(23IM30010)

\vspace{0.5in}

Under the supervision of\\
{\bfseries Prof. S. P. Sarmah}

\vspace{0.5in}

\includegraphics[width=1.35in]{logo_ready.png}

\vspace{0.5in}

Department of Industrial and Systems Engineering\\
Indian Institute of Technology Kharagpur\\[4pt]
{\bfseries Autumn Semester, 2026--27}\\
September 17, 2026

\vfill
\end{center}
\end{titlepage}

% =========================================================
% HEADER
% =========================================================
\twocolumn[
\begin{@twocolumnfalse}

\begin{center}

{\large\bfseries
Dynamic Truck--Drone Collaborative Delivery with Flexible Launch Positions
}

\vspace{5pt}

{\bfseries Jeet Anand (23IM30010)}

\vspace{2pt}

{\small
Department of Industrial and Systems Engineering,
Indian Institute of Technology Kharagpur
}

\vspace{2pt}

{\small
BTP-I Project Summary --- Autumn 2026--27
\quad$\bullet$\quad
Supervisor: Prof. S. P. Sarmah
}

\end{center}

\vspace{3pt}

\end{@twocolumnfalse}
]

% =========================================================
% ABSTRACT
% =========================================================

\textit{\textbf{Abstract}---Truck--drone collaborative delivery with
flexible launch positions has been studied as a static planning
problem and, more recently, as a dynamic problem solved through a
learned dispatch policy; separately, dynamic re-planning of an
already-executing delivery route has been studied for single-vehicle
attended-delivery systems using local-search heuristics. This project
first reproduces the static formulation of Han, Chi, and Sun [1] as a
validated baseline, then investigates a computational mechanism for
requests that arrive while a plan is already executing that differs
from both dynamic references: localized modification of the unexecuted
portion of the plan under customer-specific delay tolerances, in
contrast to the Markov-decision-process/reinforcement-learning policy
of Cui et al. [2] and the whole-route re-optimization used by
Chakraborty et al. [3]. This summary states the problem, positions it
relative to the three references, and reports current implementation
status.}

% =========================================================
\section{Motivation and Research Context}
% =========================================================

Combining a road-based truck with air-based drones lets the two
vehicles serve different customers in parallel. Han et al. [1]
formalize a static version in which a drone may launch from a flexible
point along the truck's route, not only from a customer stop, and is
recovered at a truck-service node or the ending depot. The full
customer set is known when the plan is built, and the objective is
total customer waiting time, solved by an MILP-guided VNS heuristic.

Cui et al. [2] address a genuinely dynamic truck--drone problem:
requests for a fleet of multiple trucks and drones arrive randomly,
each with its own time window and a delay penalty, and both launch and
retrieval positions are chosen en-route. They cast the real-time
decisions as a Markov decision process and solve them with a
reinforcement-learning deep Q-network policy that decides, per
request, whether and by which vehicle it is served.

A separate line of work addresses dynamic re-planning without drones.
Chakraborty, Darbhe, and Sarmah [3] study attended home delivery for
India's public distribution system with a single vehicle per depot.
A newly arriving request is checked against an already-planned day's
route and, if accepted, an iterated-local-search/adaptive-large-
neighborhood-search metaheuristic is rerun over the unexecuted route.

This project poses the dynamic question differently. Rather than
learning a policy over arrival states or re-optimizing the whole
remaining route, Stage 2 asks how an already-executing truck--drone
plan produced by Stage 1 should be locally repaired: the completed
portion is held fixed, changes are confined to a bounded region around
the new request, and previously scheduled customers are required to
remain within their individual delay tolerances.

% =========================================================
\section{Static Baseline and Dynamic Problem}
% =========================================================

Stage 1 reproduces Han et al.'s static problem: one truck, multiple
drones, a known customer set, flexible launch positions (FLPs) generated
along the arcs actually traversed by the truck, drone recovery
restricted to a truck-service node or the ending depot, drone-endurance
and truck/drone synchronization constraints, and a total-waiting-time
objective. Flexible recovery is therefore not introduced, unlike the
en-route retrieval considered by Cui et al.

Stage 2 begins once a plan $S$ produced by Stage 1 has been validated
and is executing. When a new request $c_{\mathrm{new}}$ arrives, the
executed portion of $S$ is frozen. A candidate modification $S'$ may
only act on the remainder and is confined to an affected region around
$c_{\mathrm{new}}$, such as nearby customers, the local route segment,
idle drones, and nearby FLPs. Flexible launch positions are retained;
flexible recovery is not introduced.

For each existing customer $i$, with planned delivery time $T_i$ and
modified delivery time $T_i'$, define the delivery delay as

\begin{equation}
d_i = T_i' - T_i,
\qquad
d_i \leq T_i^h ,
\label{eq:tolerance}
\end{equation}

where $T_i^h$ is the customer's allowable delay tolerance relative to
its own scheduled time.

Candidates are compared using the change in total waiting time
$\Delta W$, whether $c_{\mathrm{new}}$ is served, the number of
existing customers affected, tolerance violations, and computation
time. A working scalarization is

\begin{equation}
\min\quad
\Delta W-\alpha N_{\mathrm{new}}
+
M\sum_{i\in C_{\mathrm{old}}}v_i ,
\label{eq:objective}
\end{equation}

where

\begin{equation}
v_i=\max\left(0,d_i-T_i^h\right).
\label{eq:violation}
\end{equation}

Here $\alpha$ and $M$ are design parameters to be calibrated from
baseline results; they are not established values.

% =========================================================
\section{Proposed Mathematical Formulation}
% =========================================================

Let $S$ denote the currently executing plan and $S'$ a candidate
modified plan. The Stage-2 decision is to select a feasible local
modification of $S$ that accommodates the new request while limiting
the disruption to existing customers.

The proposed formulation is represented compactly as

\begin{equation}
\begin{aligned}
\min_{S'}\quad
&\Delta W-\alpha N_{\mathrm{new}}
 +M\sum_{i\in C_{\mathrm{old}}}v_i \\[1mm]
\text{s.t.}\quad
&S'_{\mathrm{executed}}=S_{\mathrm{executed}},\\
&S'_{\mathrm{changed}}\subseteq R(c_{\mathrm{new}}),\\
&d_i\leq T_i^h,
\qquad i\in C_{\mathrm{old}},\\
&\text{truck--drone synchronization},\\
&\text{drone endurance},\\
&\text{drone fleet capacity},\\
&\text{route and service feasibility}.
\end{aligned}
\label{eq:formulation}
\end{equation}

Here $R(c_{\mathrm{new}})$ denotes the affected region selected around
the new request. Its exact definition and size are research variables
to be investigated experimentally. The formulation is therefore a
working Stage-2 model rather than a finalized mathematical program.

The principal design distinction is that the completed portion of the
original plan is immutable, while only a bounded portion of the
unexecuted plan may be changed. Candidate plans are evaluated using
the same deterministic timing and feasibility engine developed for
Stage 1.

% =========================================================
\section{Methodology and Implementation}
% =========================================================

\textbf{Stage-1 pipeline (implemented):}
instance $\rightarrow$ truck-only TSP $\rightarrow$ $k$-means clustering
$\rightarrow$ FLP generation on traversed arcs $\rightarrow$ greedy
drone-trip construction with temporal fleet-capacity checking
$\rightarrow$ initial solution. A VNS driver with eight local-search
and four shake neighborhoods, following Han et al., is designed to
refine the solution but is not yet fully implemented.

Every stage calls one timing/feasibility engine that is the sole place
computing arrival times, synchronization, endurance, fleet capacity,
feasibility, and the objective. Thus, Stage 2 will later evaluate
candidate modifications using the same deterministic rules.

\textbf{Stage-2 pipeline (planned):}
executing plan $\rightarrow$ new request $\rightarrow$ freeze committed
portion $\rightarrow$ identify affected region $\rightarrow$ generate
local modifications (truck insertion, drone assignment, launch-position
change, local reordering) $\rightarrow$ evaluate candidates through
the Stage-1 timing engine $\rightarrow$ select a feasible modification
within tolerance, or reject.

Neighborhood-based and clustering-based affected-region definitions
are candidate strategies to be compared experimentally. Unlike the
day-level acceptance test of Chakraborty et al. [3], the decision unit
here is a bounded region rather than the entire remaining route.

% =========================================================
\section{Current Progress}
% =========================================================

\begin{center}
\footnotesize
\textbf{TABLE I.} Current Implementation Status

\vspace{2pt}

\begin{tabular}{p{0.65\columnwidth}c}
\toprule
\textbf{Component} & \textbf{Status}\\
\midrule
Core data model, geometry \& travel-time layer & \textbf{Done}\\
Flexible launch point generation & \textbf{Done}\\
Timing \& feasibility engine & \textbf{Done}\\
Initial-solution construction & \textbf{Done}\\
VNS driver and neighborhoods & \textit{In progress}\\
Experiment runner \& validation vs. Han et al. & Planned\\
Stage-2 dynamic layer & Planned\\
\bottomrule
\end{tabular}
\end{center}

The completed components are exercised by a regression suite of 339
passing assertions covering the data model, geometry, FLP generation,
timing/feasibility, and initial-solution construction. The VNS driver
and downstream components are not yet covered because they are still
under development.

% =========================================================
\section{Positioning Relative to Reference Studies}
% =========================================================

\begin{center}
\footnotesize
\textbf{TABLE II.} Feature Comparison with Reference Studies

\vspace{2pt}

\setlength{\tabcolsep}{2.5pt}
\begin{tabular}{lcccccccc}
\toprule
 & SP & DR & TD & FL & FR & RL & LR & DT\\
\midrule
Han et al. [1]
 & \checkmark & -- & \checkmark & \checkmark
 & -- & -- & -- & --\\

Cui et al. [2]
 & \checkmark & \checkmark & \checkmark & \checkmark
 & \checkmark & \checkmark & -- & --\\

Chakraborty et al. [3]
 & \checkmark & \checkmark & -- & --
 & -- & -- & WR & --\\

\textbf{This work}
 & \checkmark\,(S1) & \checkmark\,(S2) & \checkmark & \checkmark
 & -- & -- & \checkmark & \checkmark\\

\bottomrule
\end{tabular}
\end{center}

\footnotesize
\textbf{SP}: static planning; \textbf{DR}: dynamic requests;
\textbf{TD}: truck--drone coordination; \textbf{FL}: flexible launch;
\textbf{FR}: flexible recovery; \textbf{RL}: learned RL/MDP policy;
\textbf{LR}: localized repair; \textbf{WR}: whole remaining route;
\textbf{DT}: per-customer delay tolerance. For \textbf{This work},
\textbf{(S1)} and \textbf{(S2)} indicate the project stage: S1 is the
static Stage-1 baseline, already implemented and tested, while S2 is
the proposed dynamic Stage-2 extension.

% =========================================================
\section{Research Direction and Remaining Work}
% =========================================================

The immediate task is finishing the VNS driver and its twelve
neighborhoods, followed by an experiment runner that reproduces Han
et al.'s synthetic instances and checks objective values and runtime
against their reported results. Mismatches arising from unspecified
paper details will be distinguished from implementation errors.

Only after this baseline is validated will Stage 2 be implemented. The
next components are the event layer, affected-region strategies, and
local-modification generator, followed by single-request experiments.
These experiments will compare region definitions and sizes using
$\Delta W$, tolerance violations, number of affected customers, and
computation time.

The Stage-2 direction differs from Cui et al.'s trained dispatch policy
because no policy is learned across episodes. It also differs from
Chakraborty et al. [3] because acceptance and modification are confined
to a bounded region around the arriving request rather than the whole
remaining route.

% =========================================================
\section{Conclusion}
% =========================================================

Stage 1 reproduces Han et al.'s static flexible-launch problem, with
the core computational pipeline implemented and independently tested;
VNS refinement and paper-level validation remain. Stage 2 will
investigate a different dynamic mechanism: localized modification of
an already-executing truck--drone plan under customer-specific delay
tolerances. The objective is to accommodate newly arriving requests
while limiting disruption to the existing plan, using the validated
Stage-1 timing and feasibility framework as its computational
foundation.

% =========================================================
% REFERENCES
% =========================================================

\balance

\begin{thebibliography}{3}
\setlength{\itemsep}{0pt}
\setlength{\parsep}{0pt}
\footnotesize

\bibitem{han}
B. Han, T. Chi, and Z. Sun,
``Collaborative delivery problem of trucks and multiple drones:
an en-route operation with flexible launch positions,''
\textit{International Transactions in Operational Research},
vol. 33, pp. 4072--4109, 2026,
doi: 10.1111/itor.70085.

\bibitem{cui}
H. Cui, K. Li, S. Jia, and Q. Meng,
``Dynamic collaborative truck-drone delivery with en-route
synchronization and random requests,''
\textit{Transportation Research Part E: Logistics and
Transportation Review}, vol. 192, p. 103802, 2024,
doi: 10.1016/j.tre.2024.103802.

\bibitem{chakraborty}
S. Chakraborty, K. Darbhe, and S. Sarmah,
``Attended home delivery in Indian public distribution system:
an iterated local search approach,''
\textit{Journal of Modelling in Management},
vol. 16, no. 4, pp. 1279--1301, 2021,
doi: 10.1108/JM2-06-2020-0148.

\end{thebibliography}

% =========================================================
% SIGNATURES
% =========================================================

\vspace{3pt}

\begin{center}

\begin{tabular}{
>{\centering\arraybackslash}p{0.40\columnwidth}
@{\hspace{0.12\columnwidth}}
>{\centering\arraybackslash}p{0.40\columnwidth}
}

\rule{0.90\linewidth}{0.4pt}
&
\rule{0.90\linewidth}{0.4pt}
\\[3pt]

\textbf{Jeet Anand (23IM30010)}
&
\textbf{Prof.\ S.\ P.\ Sarmah}
\\[2pt]

Student
&
Supervisor
\\[5pt]

Date: \rule{0.55\linewidth}{0.4pt}
&
Date: \rule{0.55\linewidth}{0.4pt}

\end{tabular}

\end{center}

\end{document}
